% STATUS: DRAFT
\chapter{Superselection and DHR}\label{ch:dhr}
\skippable

\begin{physmotiv}
How do electric charge, baryon number and spin-statistics show up in a theory formulated only through local observables? Doplicher, Haag and Roberts showed that charged sectors are \emph{other representations} of the same quasi-local algebra (\cref{ch:inequiv}), distinguished by being indistinguishable from the vacuum outside a bounded region. The set of such sectors forms a category with products (fusion) and a braiding. Statistics (Bose/Fermi, or parastatistics, or anyons in $2+1$ dimensions) is read off from the braiding, and in $d\ge3+1$ the gauge group is recovered (Doplicher--Roberts). This chapter is a conceptual overview; it is not needed for the gravity chapters.
\end{physmotiv}

\section{DHR sectors}

Let $\pi_0$ be the vacuum representation of the quasi-local algebra $\A$. A representation $\pi$ satisfies the \emph{DHR selection criterion} if for every double cone $O$,
\begin{equation}
\pi\restriction\A(O')\ \cong\ \pi_0\restriction\A(O')\quad\text{(unitarily equivalent)}.
\end{equation}
This says the charge is localized in a bounded region: far from it, $\pi$ looks like the vacuum. Such $\pi$ can be realized as $\pi_0\circ\rho$ with $\rho$ a \emph{localized endomorphism} of $\A$ ($\rho=\id$ on $\A(O')$). Equivalence classes of irreducible such $\rho$ are the \emph{superselection sectors}.

\section{Structure}
\begin{itemize}
\item \textbf{Composition.} $\rho_1\circ\rho_2$ is again localized: charges can be fused. The sectors form a (tensor) category $\mathrm{DHR}(\A)$ with direct sums, conjugates (antiparticles, for finite statistics) and the vacuum as unit.
\item \textbf{Statistics.} Localizing $\rho_1,\rho_2$ at spacelike separation, the two compositions $\rho_1\rho_2$ and $\rho_2\rho_1$ are unitarily equivalent by a \emph{statistics operator} $\varepsilon(\rho_1,\rho_2)$. In $d\ge3+1$, $\varepsilon^2=1$ on irreducible sectors: Bose or Fermi (or parastatistics of finite order, then resolved by the next theorem). In $2+1$ dimensions $\varepsilon$ generates a braid group representation: anyons.
\item \textbf{Dimension.} Every sector has a dimension $d(\rho)\ge1$, $d(\rho)=1$ iff abelian charge; $d(\rho)$ is the \emph{quantum dimension} and determines how a sector contributes to entropy (for instance, topological entanglement entropy).
\end{itemize}

\begin{theorem}[Doplicher--Roberts, $d\ge3+1$, conceptual statement]
Under the above conditions there is a compact group $G$ and a net of ``field'' algebras $\mathcal F(O)$ with graded locality, carrying a $G$-action whose fixed points are $\A(O)$, and whose sectors are exactly the representations of $G$. The gauge group is reconstructed from the observables.
\end{theorem}

In this sense gauge symmetry is not an input, it is a derived property of a local net. It also fits with the central-decomposition viewpoint of \cref{ch:factors}: superselection sectors are labels of a decomposition (a generalized centre) of the algebra of \emph{all} (including charged) operators.

\begin{whatwrong}
DHR theory assumes a mass gap or at least Haag duality and a localization of the charge in a \emph{bounded} region; theories with long-range forces (electric charge with massless photons, Gauss's law at infinity, infraparticles) require a modified analysis (Buchholz--Fredenhagen). And, in gravity, \cref{ch:gravity_constraints} will argue there are no local charges and no non-trivial sectors localized in bounded regions.
\end{whatwrong}

\section*{Exercises}
\begin{exercise}
Show that a localized endomorphism of the Weyl algebra of a free field, $\rho(W(f))=\ee^{\ii\lambda\,\ell(f)}W(f)$ with $\ell$ a symplectic functional supported in $O$, is DHR-localized, and describe its charge (a coherent shift).
\end{exercise}
\begin{exercise}
Using the graded commutation of fermionic operators, show that the fermion parity $(-1)^F$ is a $\Z_2$ gauge group whose representations give two sectors (even/odd).
\end{exercise}
